\section{Background}
\label{sec:background}

The deployment of AI agents capable of autonomous planning, tool use, and multi-step execution has grown substantially in recent years, driven by advances in large language models (LLMs) and orchestration frameworks \cite{maeda_etal_2024_structuredagent, wu_etal_2024_magicscript}. As these agents proliferate in enterprise, scientific, and operational contexts, managing their lifecycles, enforcing accountability, and preserving provenance across delegation chains has become a central challenge in AI governance. This section situates \textbf{Recursive Spawning} within this landscape by reviewing: (1) agent lifecycle management in multi-agent systems; (2) accountability and provenance mechanisms in distributed AI; (3) the distinction between fixed and mutable delegation chains; (4) hierarchical task decomposition as a foundation for scalable multi-agent planning; and (5) the BX3 Framework as a substrate for immutable parent chain enforcement.

\subsection{Agent Lifecycle Management in Multi-Agent Systems}
\label{sec:background_lifecycle}

Agent lifecycle management encompasses the full arc from design and training through deployment, monitoring, optimization, and decommissioning \cite{onereach_2026_agent_lifecycle}. Traditional software agent frameworks treated deployment as a terminal event: an agent is invoked, executes, and returns a result. Modern agentic systems, by contrast, treat agents as long-lived entities that must adapt to evolving contexts, recover from failures, coordinate with other agents, and remain aligned with human objectives over extended operational horizons \cite{weng_2025_selfmodification, beebe_2026_bx3_grant}. This shift introduces significant complexity in governance. The lifecycle management literature distinguishes multiple phases---design, training, testing, deployment, monitoring, optimization, and decommissioning---each with distinct accountability requirements \cite{onereach_2026_agent_lifecycle, SPIE_2026_assurance}. In hierarchical multi-agent architectures, these phases become nested: a parent orchestrator agent may spawn, configure, monitor, and retire subordinate agents as sub-tasks complete or conditions change. Without systematic lifecycle management, agents can accumulate undocumented behavioral drift, creating accountability gaps that are difficult to audit retroactively.

Frameworks such as the Hierarchical Task DAG (HTDAG) introduced by Deep Agent demonstrate how dynamic, DAG-based decomposition enables real-time reorganization of sub-tasks as objectives evolve \cite{zhang_etal_2025_deepagent}. Similarly, the Agent Control Protocol (ACP) provides deterministic, history-aware admission control over agent actions in distributed environments, enforcing hard boundaries rather than advisory constraints \cite{ACP_2025}. These systems illustrate that effective lifecycle management requires both architectural discipline (hierarchical decomposition, isolation of concerns) and enforcement mechanisms (admission control, audit trails) that persist across agent spawn, execution, and retirement events.

\subsection{Accountability and Provenance in Distributed Systems}
\label{sec:background_provenance}

Distributed AI systems face a fundamental accountability problem: when a multi-agent chain produces an outcome, determining which actor---human or machine---authorized, executed, or contributed to each step is non-trivial. This problem is compounded when agents operate across trust boundaries, delegate authority to one another, or execute in adversarial environments. Provenance tracking addresses this by maintaining a verifiable record of action authorship, delegation scope, and authorization chain \cite{haberkamp_2025_ipp, beyer_etal_2025_agent_identity, helixar_2026_hdp}.

Several emerging standards tackle this challenge. The Intent Provenance Protocol (IPP) defines a cryptographic mechanism to carry and verify human intent as agents act, using Ed25519 signatures and Decentralized Identifiers (DIDs) to create tamper-evident, scope-bounded intent tokens \cite{haberkamp_2025_ipp}. The Human Delegation Provenance (HDP) protocol extends this model to hierarchical multi-agent deployments, creating an append-only, cryptographically signed delegation chain in which each hop records the delegator, delegatee, authorization scope, and session context \cite{helixar_2026_hdp, dibbo_etal_2026_hdp_formal}. HDP enables offline verification of the full provenance chain using only the issuer's public key, without registry lookups or trusted third parties. The Agent Identity Architecture draft proposes human-anchored identity roots, explicit delegation semantics, and provenance structures that trace actions back to responsible humans across contexts and platforms \cite{beyer_etal_2025_agent_identity}.

These protocols share a common recognition: accountability is not a property that can be retroactively injected into a system. It must be architecturally baked into the delegation chain from inception. The SentinelAgent framework formalizes this through its Delegation Chain Calculus (DCC), which specifies seven properties including \emph{authority narrowing}, \emph{policy preservation}, and \emph{forensic reconstructibility} \cite{sentinel_2025}. DCC's seven properties guarantee that delegated authority narrows at each hop, that policies are preserved across delegation levels, and that any auditor can reconstruct the full delegation history. Experimental evaluation against 516 adversarial scenarios across 13 federal domains achieved a 100\% true positive rate at 0\% false positive rate, demonstrating that formal delegation semantics can enforce accountability at scale \cite{sentinel_2025}.

Blockchain-based approaches offer an alternative or complement to cryptographic delegation tokens. GuardianChain's GC-1 system anchors AI decision capsules to public blockchains via hash-linked seals, enabling any party to verify the integrity of a decision chain without trusting a central authority \cite{cronin_2025_guardianchain}. Similarly, implicit execution tracing (IET) embeds agent-specific signals into token distributions during generation, enabling reconstruction of the delegation and coordination graph from final output alone, while preserving privacy \cite{iet_2025}. These approaches are particularly relevant for regulated environments (healthcare, defense, finance) where auditability is a compliance requirement, not merely a best practice.

\subsection{Fixed vs. Mutable Delegation Chains}
\label{sec:background_delegation}

A critical design decision in any delegation architecture is whether the delegation chain itself is mutable or immutable after instantiation. In a \textbf{mutable delegation chain}, parent pointers, authorization scopes, and capability bindings can be modified post-deployment---either by the delegator, by an administrator, or by some other authorized actor. Mutable chains offer flexibility: scopes can be widened as trust is established, agents can be re-authorized for new tasks, and administrative corrections can be applied without redeployment. However, mutability introduces a fundamental accountability hazard: the chain of authorization that produced a given action can be altered after the fact, undermining the evidentiary value of any audit trail \cite{cronin_2025_guardianchain}.

In contrast, a \textbf{fixed (immutable) delegation chain} establishes parent-child relationships, authorization scopes, and capability bindings at spawn time and prevents subsequent modification. Once a child agent is instantiated with a given Purpose pointer and bounded scope, those bindings are cryptographically or architecturally sealed. Any subsequent attempt to widen scope, redirect the parent pointer, or modify capability bindings constitutes a new delegation event, requiring fresh authorization and generating a new chain entry.

The distinction has practical consequences for autonomous drift. When a child agent operates with a mutable parent pointer, it can---in principle---be re-parented to a different Purpose layer without its knowledge, redirecting its reasoning and execution toward an unauthorized objective. This failure mode is especially dangerous in edge deployments where cloud connectivity is intermittent and agents must operate independently for extended periods. The literature on agentic AI delegation identifies scope creep and unauthorized re-delegation as primary attack vectors in multi-agent systems \cite{dibbo_etal_2026_hdp_formal, wu_etal_2024_governance}. Immutable parent chains structurally prevent this: the child agent's Purpose pointer is fixed at instantiation, and any change requires a new spawn event with a new cryptographic record, preserving full auditability.

The distinction between mutable and fixed delegation is reflected in existing standards. HDP's append-only chain inherently favors immutability: each hop is signed and additions do not alter prior entries \cite{helixar_2026_hdp}. IPP's \emph{Narrowing Invariant} enforces that derived tokens never exceed their parent's authorized scope, effectively treating scope as an immutable property of the delegation instance \cite{haberkamp_2025_ipp}. SentinelAgent's \emph{authority narrowing} property formally guarantees that each delegation hop reduces or preserves authority, never expanding it \cite{sentinel_2025}. These mechanisms collectively establish that fixed delegation chains are not merely a design preference but an emerging consensus in accountable agent architecture.

\subsection{Hierarchical Task Decomposition in AI}
\label{sec:background_htd}

Hierarchical task decomposition is the process of recursively breaking high-level objectives into semantically meaningful subgoals, each assigned to a specialized agent or module. This approach addresses the fundamental scaling limitation of flat agent architectures: a single monolithic agent reasoning over long-horizon, multi-constraint objectives suffers from in-context memory saturation, error propagation across steps, and loss of task structure \cite{mao_etal_2025_hiplan, chen_etal_2025_reactree, lu_etal_2025_hipere}.

The HiPlan framework introduces a two-layer decomposition separating global milestones (macroscopic guidance) from local steps (real-time adaptive actions), enabling agents to preserve task structure while remaining responsive to environmental changes \cite{mao_etal_2025_hiplan}. ReAcTree extends this with dynamic agent tree construction where each node is an LLM-powered agent capable of reasoning, acting, and expanding its subtree, achieving 61\% goal success on long-horizon benchmarks versus 31\% for flat ReAct baselines \cite{chen_etal_2025_reactree}. STRUCTUREDAGENT uses AND/OR tree planning with structured memory to maintain candidate solutions across extended task horizons, demonstrating improved interpretability and human intervention capability \cite{maeda_etal_2024_structuredagent}. HiMAC bi-level hierarchy separates macro-level planning (planner) from micro-level execution (executor), reducing exploration complexity and mitigating error cascades in long-horizon tasks \cite{lu_etal_2025_hipere}.

Mechanistically interpretable task decomposition (MITD) provides diagnostic transparency by splitting reasoning into Planner, Coordinator, and Executor modules with attention-based visualizations, showing that optimal decomposition depths of 12--25 steps reduce reward hacking frequency by approximately 34\% \cite{mitd_2025}. MagicAgent addresses data scarcity in hierarchical planning through synthetic trajectory generation, demonstrating that structured decomposition can be learned and generalized across diverse task domains \cite{wu_etal_2024_magicscript}. The Intelligent AI Delegation framework extends hierarchical decomposition to encompass not just task breaking but full accountability semantics: who delegates to whom, under what authority, with what monitoring and recovery mechanisms \cite{zhang_wei_2025_intelligent_ai_delegation}.

These works collectively establish that hierarchical decomposition is essential for scalable, accountable multi-agent systems. However, they also reveal a gap: most hierarchical frameworks focus on the \emph{planning and execution} dimensions of decomposition but provide limited architectural support for the \emph{accountability and provenance} dimension as agents spawn, delegate, and retire across levels. Recursive Spawning addresses this gap by ensuring that every spawned child inherits a sealed, immutable parent pointer, making the delegation chain itself a first-class architectural primitive.

\subsection{The BX3 Framework as Substrate for Immutable Parent Chains}
\label{sec:background_bx3}

The BX3 Framework (Universal Architecture for Accountable Autonomous Systems) defines a three-layer functional architecture---Purpose, Bounds Engine, and Fact Layer---that decouples the \emph{why} of an action (Purpose, Human-anchored) from the \emph{how} (Bounds Engine, limbless proposer) and the \emph{what} (Fact Layer, deterministic physical executor) \cite{beebe_2026_bx3_zenodo, beebe_2026_bx3_grant}. This separation ensures that human accountability is structurally preserved regardless of which actor occupies each layer, and that no reasoning entity shares a functional plane with its physical executor \cite{beebe_2026_bailout}.

Pillar 2 of the BX3 Framework, \textbf{Recursive Spawning}, operationalizes hierarchical decomposition through a controlled spawn mechanism. A Parent Node births a Child Loop by provisioning a \textbf{Worksheet}: a self-contained, containerized logic set delivered Over-the-Air (OTA) to the target node. The Worksheet contains the complete Bounds Engine and Fact Layer logic for the target environment, enabling the child to execute last-known-good reasoning during cloud disconnection (Local Survivability). Critically, each child retains a \textbf{hard-coded pointer to its Parent's Purpose layer}, preventing autonomous drift even when the child operates independently of the parent for extended periods.

The immutability of the parent pointer is architecturally enforced at the Fact Layer: the pointer is a deterministic, physical constraint expressed as a SHA-256 hash of the parent's canonical Purpose, sealed at spawn time. Any subsequent modification to the parent pointer constitutes a new spawn event, generating a fresh cryptographic record and severing the continuity of the delegation chain. This design closes the mutable delegation vulnerability identified in Section \ref{sec:background_delegation} by making parent chain integrity a physical-layer invariant, not a policy-layer convention.

The BX3 Framework's Forensic Ledger complements immutable parent chains with an append-only, hash-linked event record that captures every state transition, every spawn event, and every bounds check across the full agent hierarchy \cite{beebe_2026_bx3_zenodo}. Each entry is sealed with the generating node's Fact Layer hash, ensuring that any tampering with historical records is detectable. The combination of immutable parent pointers (structural constraint) and forensic sealing (evidentiary record) provides a dual guarantee: not only can delegation chains not be silently redirected, but any attempt to do so leaves a cryptographically verifiable trace.

Formal analysis of the BX3 Framework's accountability properties references the Delegation Chain Calculus (DCC) \cite{sentinel_2025} as a comparative baseline. The BX3 Purpose Layer maps directly to DCC's human-anchored authority root; the BX3 Bounds Engine enforces DCC's scope-action conformance property; and the BX3 Fact Layer implements DCC's forensic reconstructibility guarantee through deterministic physical enforcement. Where DCC provides the formal specification, the BX3 Framework provides the concrete implementation substrate---making abstract delegation semantics executable in real-world edge deployments.

\subsection{Summary}
\label{sec:background_summary}

The background literature reveals a field in active convergence. Agent lifecycle management frameworks have established the necessity of treating agents as long-lived, governed entities rather than stateless functions. Accountability and provenance protocols have demonstrated that cryptographically verifiable delegation chains are technically achievable and practically deployable. The fixed-versus-mutable delegation distinction has sharpened into a consensus that immutable parent pointers are a prerequisite for accountable autonomous operation. Hierarchical task decomposition has proven essential for scaling multi-agent planning, but its integration with accountability semantics remains an open frontier. The BX3 Framework's Recursive Spawning mechanism occupies this frontier by providing: (1) a structured spawn mechanism that decomposes tasks hierarchically; (2) immutable parent Purpose pointers that prevent autonomous drift; (3) a forensic ledger that creates verifiable, tamper-evident spawn and execution records; and (4) a physical-layer enforcement model that makes these guarantees deterministic rather than advisory.

% Bibliography
% Note: compiled separately with bibtex using refs.bib
